% Artículo VII — primera redacción editorial, 10 de septiembre de 2026.
% Apertura integrada con las demostraciones de la revisión 06.
% Procedencia y condiciones: PLAN_RESIDENCIAS_DEMOSTRATIVAS.md.
% Las referencias internas remiten a demostraciones incluidas en el cuerpo.

\begin{center}\bfseries Résumé\end{center}
\begingroup
\fontsize{9.5}{10.7}\selectfont
\setlength{\parskip}{1.5pt plus .5pt minus .5pt}
\noindent
L'amplitude torsionnelle est déterminée comme fonctionnelle explicite de l'état matériel au moyen de l'action spinorielle, de l'élimination de la contorsion et de la variation métrique. On construit un secteur matériel dans lequel cette amplitude est positive et se conserve. Sa contribution énergétique produit une dilution \(a^{-6}\) et, dans la réalisation homogène plate avec \(\Lambda_{\mathrm{eff}}=0\), des densités positives et \(w<1\), un seuil de rebond unique et transversal. Pour la poussière, on obtient une solution exacte pour tout temps, à courbure finie et à géodésiques causales de Levi--Civita complètes dans les deux directions.

La construction part d'un noyau formel nommé holographie modulaire
triadique, ci-après HMT. \emph{All Possible Paths}, ci-après APP,
réunit des chemins arithmétiques sur \((\mathbb Z/9\mathbb Z)^2\), avec une
adjacence donnée par le graphe de Cayley additif de générateurs cardinaux
et une réalisation cellulaire toroïdale, en distinguant les opérations et les
évaluations de somme et de produit ;
TRIT est l'unité de signature orientée \(\{-1,0,+1\}\) qui détermine les
régimes locaux ; le \emph{Topological Phase Kernel}, ci-après TPK,
compose la dynamique de transport et de mémoire. L'état enrichi
conserve les feuillets, l'orientation, le résidu, la retenue et l'histoire pendant le
retour de phase. La mécanique dimensionnelle exprime ses résultats
internes au moyen de grandeurs et d'observables physiques, en préservant la
provenance structurelle des couplages.

La normalisation gravitationnelle est développée au préalable au moyen de la
composition longitudinale \(L=(5759/23040)\alpha^{16}L_*\) et de sa
métrique radiale positive. Dans la réalisation définie par ces opérations
et par l'identification du rayon gravitationnel réduit, la réciprocité
avec la section d'action fixe \(G=c^3L^2/\hbar_{\rm ret}\) dans tous
les modes admissibles. La démonstration conserve l'archive complète et
établit les conditions de réduction statique et dynamique de la mémoire.
Le choix dimensionnel \(L_*=1\,\mathrm m\), les règles longitudinale
et métrique et l'identification radiale sont déclarés dans l'énoncé même.

Le transport des routes admet une réalisation de Cartan compatible avec la composition, l'inversion et le raffinement. Sur celle-ci, on distingue deux actions matérielles : l'extension minimisante de l'écran de mémoire et la réalisation spinorielle. Chacune conserve sa propre variation par rapport à la connexion. Dans le secteur spinoriel, affine en la contorsion, l'inversion des opérateurs de Holst et de torsion détermine la correction effective ; pour l'extension minimale, on conserve la dépendance non linéaire du courant et on démontre un critère local au moyen de la hessienne. Le courant, l'observable quartique et sa normalisation expliquent la valeur de l'amplitude, outre son rôle dans l'évolution cosmologique.

La persistance du rebond sous perturbations se démontre au moyen de bornes locales différenciées, tandis que la complétude causale correspond à la solution globale explicite de poussière. Les bilans résiduels, l'information d'horizon et les opérations de l'observateur conservent les conditions relatives aux sources, aux couplages et aux coordinatisations qu'ils utilisent. La réponse relationnelle de l'état conjoint se relie à la propagation avancée et retardée au moyen d'opérateurs de Green pour des sources à support fini, associés à la différence covariante construite à partir du transport des parcours et d'une condition de compatibilité à la frontière. Le résultat relie la mémoire engendrée, sa réalisation géométrique et les conséquences dynamiques et informationnelles de cette réalisation.

La composition entre action, incrément surfacique de Barbero--Immirzi et
horizon détermine la conversion
\(T_{\rm cos}(L_\alpha)/\Theta_{\rm clk}=\log3/(2\pi)\).
L'hamiltonien de réponse angulaire fournit un bilan fini
d'énergie libre au moyen de l'entropie relative et une première loi à rayon
fixe. Une section thermométrique relative caractérise en outre le
transducteur positif et le critère permettant d'identifier un autre lecteur sur
cette même section.

\noindent\textbf{Mots-clés :} gravitation ; torsion ; mémoire arithmétique ; équations de champ ; longueur de Planck ; transition de la contraction à l'expansion.
\par\endgroup
